\documentclass[a4paper,12pt,oneside,openany]{book}

% ============================================================================
% PROJETO FINAL - BANCO DE DADOS I
% UNESP - Faculdade de Ciencias - Bauru - Departamento de Computacao
% Prof. Bruno Elias Penteado
%
% Projeto: NextMedia - plataforma de streaming de video por assinatura
%
% ENTREGAVEIS PEDIDOS PELO PROFESSOR (orientacao complementar):
%   Grupo de 4 pessoas | cenario escolhido pelo grupo | 10 a 15 entidades
%   Documento PDF contendo:
%   1) Modelagem conceitual: descricao BREVE do minimundo, regras de negocio,
%      diagrama DER
%   2) Modelagem logica: aplicando conceitos de normalizacao, DER logico
%   3) Modelagem fisica: codigo SQL para PostgreSQL, contendo comandos DDL e
%      DML para criacao do banco e carga de dados iniciais de exemplo,
%      + 3 SELECTs que exemplifiquem a utilidade no minimundo
%
% COMO USAR:
% 1) Substituam os blocos \orientacao{...} pelo texto final.
% 2) Todo diagrama precisa ser citado e explicado no corpo do texto.
% 3) Antes da entrega, troquem \orientacoestrue por \orientacoesfalse.
% ============================================================================

\usepackage[T1]{fontenc}
\usepackage[utf8]{inputenc}
\usepackage[brazil]{babel}
\usepackage{graphicx}
\usepackage{booktabs}
\usepackage{tabularx}
\usepackage{longtable}
\usepackage{array}
\usepackage{enumitem}
\usepackage{xcolor}
\usepackage{listings}
\usepackage{hyperref}
\usepackage{natbib}
\usepackage{fancyhdr}
\usepackage{needspace}
\usepackage{indentfirst}
\usepackage[top=2cm,bottom=2cm,left=2cm,right=2cm]{geometry}

\hypersetup{
    hidelinks,
    pdftitle={NextMedia -- Projeto de Banco de Dados},
    pdfauthor={Equipe NextMedia}
}

\urlstyle{same}

% ----------------------------------------------------------------------------
% Listagens de codigo SQL
% ----------------------------------------------------------------------------
% Usamos o pacote "listings". NAO carregar "minted" junto: minted exige
% compilacao com -shell-escape e Pygments instalado, o que quebra o build
% para quem usa Overleaf/TeXstudio sem configuracao extra.
\definecolor{corComentario}{rgb}{0.35,0.45,0.35}
\definecolor{corPalavra}{rgb}{0.10,0.20,0.55}
\definecolor{corTexto}{rgb}{0.55,0.20,0.15}
\definecolor{corFundo}{rgb}{0.97,0.97,0.97}

\lstdefinestyle{sqlstyle}{
    language=SQL,
    morekeywords={SCHEMA,SERIAL,BIGSERIAL,TIMESTAMP,BOOLEAN,TEXT,VARCHAR,
                  NUMERIC,INTEGER,SMALLINT,DATE,INTERVAL,REFERENCES,CASCADE,
                  RESTRICT,DEFAULT,CONSTRAINT,UNIQUE,CHECK,PRIMARY,FOREIGN,
                  KEY,NOT,NULL,ON,DELETE,UPDATE,TRIGGER,FUNCTION,RETURNS,
                  LANGUAGE,PLPGSQL,BEGIN,END,IF,THEN,RAISE,EXCEPTION},
    basicstyle=\ttfamily\footnotesize,
    keywordstyle=\color{corPalavra}\bfseries,
    commentstyle=\color{corComentario}\itshape,
    stringstyle=\color{corTexto},
    backgroundcolor=\color{corFundo},
    frame=single,
    rulecolor=\color{gray},
    numbers=left,
    numberstyle=\tiny\color{gray},
    numbersep=8pt,
    breaklines=true,
    showstringspaces=false,
    tabsize=2,
    captionpos=b,
    extendedchars=true,
    inputencoding=utf8,
    literate=%
      {á}{{\'a}}1 {é}{{\'e}}1 {í}{{\'i}}1 {ó}{{\'o}}1 {ú}{{\'u}}1
      {â}{{\^a}}1 {ê}{{\^e}}1 {ô}{{\^o}}1 {ã}{{\~a}}1 {õ}{{\~o}}1
      {ç}{{\c c}}1 {Á}{{\'A}}1 {É}{{\'E}}1 {Í}{{\'I}}1 {Ó}{{\'O}}1
      {Ú}{{\'U}}1 {Â}{{\^A}}1 {Ê}{{\^E}}1 {Ô}{{\^O}}1 {Ã}{{\~A}}1
      {Õ}{{\~O}}1 {Ç}{{\c C}}1
}
\lstset{style=sqlstyle}
\renewcommand{\lstlistingname}{Listagem}
\renewcommand{\lstlistlistingname}{Lista de listagens}

% Atalho para nomes de tabelas e colunas
\newcommand{\bd}[1]{\texttt{#1}}

% ----------------------------------------------------------------------------
% Exibicao das orientacoes do modelo
% ----------------------------------------------------------------------------
% Troquem para \orientacoesfalse antes de gerar a versao final para entrega.
\newif\iforientacoes
\orientacoestrue

\newcommand{\orientacao}[1]{%
    \iforientacoes
        \begin{quote}
        \small\itshape\textbf{A escrever:} #1
        \end{quote}
    \fi
}

% Marca quem ficou responsavel pela secao (some junto com as orientacoes)
\newcommand{\autoria}[1]{%
    \iforientacoes
        \marginpar{\scriptsize\textsf{#1}}%
    \fi
}

% ----------------------------------------------------------------------------
% Capitulos em fluxo continuo
% ----------------------------------------------------------------------------
% A classe "book" comeca cada capitulo numa pagina nova, o que deixa muita
% pagina pela metade num documento curto como este. As redefinicoes abaixo
% removem essa quebra: o capitulo segue logo apos o texto anterior. O
% \Needspace evita que um titulo fique sozinho no pe da pagina.
% Para voltar ao comportamento padrao da classe, basta apagar este bloco.
\makeatletter
\renewcommand\chapter{%
    \par
    \Needspace*{5\baselineskip}%
    \thispagestyle{fancy}%
    \global\@topnum\z@
    \@afterindentfalse
    \secdef\@chapter\@schapter}
\renewcommand\@makechapterhead[1]{%
    \vspace*{1.2\baselineskip}%
    {\parindent\z@ \raggedright \normalfont
     \interlinepenalty\@M
     \large\bfseries \@chapapp\space \thechapter
     \par\nobreak
     \vskip 6\p@
     \LARGE\bfseries #1\par\nobreak
     \vskip 16\p@
    }}
\renewcommand\@makeschapterhead[1]{%
    \vspace*{1.2\baselineskip}%
    {\parindent\z@ \raggedright \normalfont
     \interlinepenalty\@M
     \LARGE\bfseries #1\par\nobreak
     \vskip 16\p@
    }}
\makeatother

% ----------------------------------------------------------------------------
% Cabecalho e rodape
% ----------------------------------------------------------------------------
\renewcommand{\headrulewidth}{0pt}
\fancyhf{}

% ----------------------------------------------------------------------------
% Informacoes do trabalho -- ATUALIZAR NOMES E RAs
% ----------------------------------------------------------------------------
\newcommand{\school}{Universidade Estadual Paulista ``Júlio de Mesquita Filho'' (UNESP)\\
Faculdade de Ciências -- Bauru\\
Departamento de Computação}
\newcommand{\discipline}{Banco de Dados I}
\newcommand{\professorname}{Prof. Dr. Bruno Elias Penteado}
\newcommand{\projtitle}{NextMedia}
\newcommand{\projsubtitle}{Projeto de banco de dados para uma plataforma de \textit{streaming} de vídeo por assinatura}
\newcommand{\semester}{2026}

% ----------------------------------------------------------------------------
% Paragrafos
% ----------------------------------------------------------------------------
% Recuo de 1,25 cm na primeira linha, como manda o uso brasileiro (e a ABNT).
% O pacote indentfirst, carregado acima, garante que o primeiro paragrafo
% depois de um titulo tambem seja recuado -- por padrao o LaTeX nao recua.
\setlength{\parindent}{1.25cm}
\setlength{\parskip}{0pt}

\graphicspath{{images/}}

\title{NextMedia -- Projeto de Banco de Dados}
\author{Equipe NextMedia}

\begin{document}

% ============================================================================
% CAPA
% ============================================================================
\thispagestyle{empty}
\begin{titlepage}
	\centering

	% Logotipo da UNESP no topo da página
	\vspace*{-12mm}
	\includegraphics[width=75mm]{AV01A.jpg}\\[8mm]

	{\large\textbf{\school}}

	\vfill

	{\huge\textbf{\projtitle}}\\[5mm]
	{\Large\projsubtitle}

	\vfill

	{\Large\textbf{\discipline}}\\[4mm]
	{\large\professorname}\\[14mm]

	%{\large\textbf{Estudantes}}\\[4mm]
	Caio César Souza Oliveira -- RA: 241025702 \\[2mm]
	Mário Masao Mukai -- RA: 241022321 \\[2mm]
	Murilo Tomaz Gonzaga -- RA: 241024684 \\[2mm]

	\vfill

	{\large Bauru -- SP\\\semester}
\end{titlepage}

\tableofcontents
\thispagestyle{empty}

\clearpage
\pagestyle{fancy}
\cfoot{\thepage}
\setcounter{page}{1}

% ============================================================================
\chapter{Introdução}\label{cap:introducao}
% ============================================================================

\section{Contexto}

A NextMedia é uma plataforma de \textit{streaming} de vídeo por assinatura.
O cliente cria uma conta, escolhe um plano e passa a assistir a um catálogo de
filmes e séries, em tantos dispositivos quantos o plano permitir. É um serviço
que a maioria das pessoas já usou, e foi por isso que o escolhemos: não é
preciso gastar páginas explicando o negócio antes de começar a modelá-lo.

Pesou também o fato de o domínio pedir, sem artifício, boa parte dos recursos
do modelo entidade-relacionamento. Um título é sempre um filme ou uma série,
nunca os dois, o que produz uma especialização total e exclusiva. Temporadas
e episódios não existem soltos, fora da série a que pertencem, e portanto são
entidades fracas. Gêneros se organizam em hierarquia e um título pode ser continuação de
outro, dois auto-relacionamentos. A participação de uma pessoa em um título,
seja como ator, diretor ou roteirista, é um relacionamento N:N que carrega
atributos próprios. Há ainda o histórico da plataforma, feito de assinaturas
encerradas e visualizações antigas, que traz o problema de preservar dados
corretos mesmo quando o catálogo e os preços mudam.

\section{Objetivo e escopo}

O objetivo do trabalho é projetar o banco de dados dessa plataforma passando
pelos três níveis de abstração usuais \citep{elmasri2019}: o conceitual, que
independe de tecnologia; o lógico, já traduzido para o modelo relacional e
normalizado; e o físico, que é um script de criação do esquema.
Usamos o PostgreSQL como sistema gerenciador, e o DER lógico foi
construído no ERD Tool, a ferramenta CASE integrada ao pgAdmin.
% TODO: completar com a ferramenta usada no DER assim que o grupo decidir.

O escopo cobre contas, planos e assinaturas, os perfis que convivem dentro de
uma mesma conta, o catálogo de títulos com suas temporadas, episódios e
gêneros, o elenco e a equipe técnica de cada obra, e o histórico de
visualizações de cada perfil. Ficam de fora o processamento de pagamentos, a
infraestrutura de entrega de vídeo, os algoritmos de recomendação e a camada de
aplicação, que pertencem à plataforma real mas não ao problema de projeto de
esquema que a disciplina propõe.

\section{Organização do documento}

Os três capítulos seguintes correspondem aos três níveis. O
Capítulo~\ref{cap:conceitual} descreve o minimundo, enuncia as regras de
negócio e apresenta o diagrama entidade-relacionamento. O
Capítulo~\ref{cap:logico} traz o modelo lógico normalizado e justifica as
decisões de mapeamento que o diagrama sozinho não explica. O
Capítulo~\ref{cap:fisico} trata da implementação no PostgreSQL e mostra qual
mecanismo garante cada regra de negócio.

% TODO: registrar aqui, em uma frase, a divisão de autoria entre os integrantes.

% ============================================================================
\chapter{Modelo Conceitual}\label{cap:conceitual}
% ============================================================================

\section{O minimundo}\label{sec:minimundo}

% ---------------------------------------------------------------------------
% ESBOCO -- primeira versao, a revisar pelo grupo.
% Checklist da secao: (i) conta, plano e assinatura, incluindo a troca de
% plano; (ii) perfis e a restricao do perfil infantil; (iii) catalogo, com
% filmes, series, temporadas, episodios, generos e franquias; (iv) pessoas e
% participacoes; (v) visualizacoes. O texto deve falar a lingua do negocio,
% deixando o vocabulario de modelagem para as secoes seguintes.
% ---------------------------------------------------------------------------

A NextMedia vende acesso a um catálogo de vídeo. O cliente abre uma conta,
identificada por um \textit{e-mail} único na plataforma, e contrata um plano,
que define preço mensal, telas simultâneas, número de perfis e qualidade. O
vínculo entre os dois é a assinatura, que guarda o período e o valor cobrado
--- registrado ali porque os planos mudam de preço. Na troca de plano, a
assinatura anterior é encerrada com data de término e uma nova é aberta, e só
uma vale por vez. Dentro da conta convivem vários perfis, um por pessoa da
casa, cada um com nome, avatar, idioma e histórico próprios; um perfil
infantil só enxerga títulos de classificação até doze anos, e o número de
perfis esbarra no limite do plano.

A unidade do catálogo é o título, com nome, ano, sinopse, idioma original e
classificação indicativa. Todo título é um filme \emph{ou} uma série: o filme
tem duração e se assiste de uma vez; a série desdobra-se em temporadas, e cada
temporada em episódios, que não existem fora da obra a que pertencem. Todo
título tem ao menos um gênero, e os gêneros formam hierarquia, já que ficção
científica é recorte de uma categoria maior. Um título também pode apontar
para o que lhe deu origem, ligando continuações e franquias.

Atores, diretores e roteiristas são registrados uma única vez, com nome e data
de nascimento. Como a mesma pessoa pode dirigir e atuar no mesmo título, o que
se registra é a participação: quem, em que obra, fazendo o quê, interpretando
qual personagem e em que posição nos créditos. Já o histórico de consumo é o
que a plataforma mais acumula: cada \textit{play} guarda qual perfil assistiu,
o quê, quando, por quanto tempo, em que dispositivo e se chegou ao fim,
apontando sempre para um filme ou um episódio, nunca para ambos. Daí também
duas decisões de retenção: um título que sai do catálogo recebe data de saída
e permanece no banco, para que as visualizações antigas continuem fazendo
sentido; a conta encerrada leva junto perfis e histórico.

O relato se traduz em treze entidades, agrupadas em quatro núcleos ---
assinatura, catálogo, pessoas e consumo:

\begin{itemize}[label=--, leftmargin=*, itemsep=2pt]
  \item \textbf{conta} --- assinante da plataforma; dados de acesso e cobrança.
  \item \textbf{plano} --- catálogo de planos comercializados.
  \item \textbf{assinatura} --- vínculo entre conta e plano, com período de vigência e valor cobrado.
  \item \textbf{perfil} --- subusuário dentro de uma conta.
  \item \textbf{titulo} --- unidade do catálogo; aponta para o título que lhe deu origem.
  \item \textbf{filme} --- título que se assiste de uma vez, com duração própria.
  \item \textbf{serie} --- título que se desdobra em temporadas.
  \item \textbf{temporada} --- depende de uma série para existir.
  \item \textbf{episodio} --- depende de uma temporada para existir.
  \item \textbf{genero} --- categoria do catálogo, podendo ter um gênero que a contém.
  \item \textbf{pessoa} --- elenco e equipe técnica.
  \item \textbf{participacao} --- atuação de uma pessoa em um título, com papel e personagem.
  \item \textbf{visualizacao} --- registro de que um perfil assistiu a um filme ou episódio.
\end{itemize}

A associação entre títulos e gêneros não entra nessa contagem: vira tabela
própria no modelo lógico, já que um título tem vários gêneros e um gênero reúne
vários títulos, mas é só o relacionamento entre duas entidades já listadas, sem
informação própria para guardar. Os atributos de cada entidade aparecem no
diagrama da Seção~\ref{sec:der} e no modelo lógico do
Capítulo~\ref{cap:logico}.

\section{Regras de negócio}\label{sec:regras}

% ---------------------------------------------------------------------------
% ESBOCO -- primeira versao, a revisar pelo grupo.
% ---------------------------------------------------------------------------

As regras a seguir saem da leitura do minimundo descrito na
Seção~\ref{sec:minimundo}. Elas justificam escolhas do diagrama que de outro
modo pareceriam arbitrárias e, principalmente, funcionam como lista de
verificação para o modelo físico: no Capítulo~\ref{cap:fisico} cada uma
reaparece ao lado do mecanismo que a faz valer. A numeração é fixa e usada como referência ao longo do
texto.\footnote{A sequência salta de RN13 para RN17 porque as regras RN14 a
RN16 tratavam de avaliações e listas pessoais, funcionalidades que ficaram
fora do escopo. Preferimos preservar a numeração das demais a renumerar tudo.}

\begin{longtable}{@{}p{0.09\textwidth} p{0.85\textwidth}@{}}
\caption{Regras de negócio da NextMedia.}\label{tab:regras}\\
\toprule
\textbf{Código} & \textbf{Regra} \\
\midrule
\endfirsthead
\multicolumn{2}{@{}l}{\small\itshape Tabela \thetable{} -- continuação}\\
\toprule
\textbf{Código} & \textbf{Regra} \\
\midrule
\endhead
\midrule
\multicolumn{2}{r@{}}{\small\itshape continua na próxima página}\\
\endfoot
\bottomrule
\endlastfoot
RN01 & O endereço de \textit{e-mail} de uma conta é único em toda a plataforma. \\
RN02 & Uma conta possui, no máximo, uma assinatura ativa por vez. \\
RN03 & O número de perfis de uma conta não excede o \bd{max\_perfis} do plano vigente. \\
RN04 & Um perfil pertence a exatamente uma conta e não pode ser transferido. \\
RN05 & Um perfil infantil só acessa títulos com classificação indicativa menor ou igual a 12 anos. \\
RN06 & Todo título é um filme \textbf{ou} uma série: a especialização é total e exclusiva. \\
RN07 & Uma temporada só existe vinculada a uma série; um episódio, apenas a uma temporada. \\
RN08 & O número da temporada é único dentro da série; o número do episódio é único dentro da temporada. \\
RN09 & Todo título possui ao menos um gênero associado. \\
RN10 & Um gênero pode ter um gênero-pai; a hierarquia não admite ciclos. \\
RN11 & A combinação (título, pessoa, papel) é única. \\
RN12 & Uma visualização refere-se a um filme \textbf{ou} a um episódio, nunca a ambos. \\
RN13 & O tempo assistido não excede a duração do conteúdo visualizado. \\
RN17 & A exclusão de uma conta propaga-se em cascata para seus perfis e seu histórico. \\
RN18 & Um título retirado do catálogo recebe data de saída e não é removido fisicamente. \\
\end{longtable}

Nem toda regra custa o mesmo para ser garantida. Algumas o próprio desenho do
modelo torna impossíveis de violar: RN04 e RN07 dispensam verificação, porque
um perfil só existe pendurado numa conta e uma temporada numa série. Outras se
resolvem com uma restrição declarada na criação da tabela --- RN01 e RN08 são
unicidade, RN11 vira a chave primária composta de \bd{participacao}, RN12 e
RN13 são verificações sobre a própria linha, e RN17 é a ação de propagação da
chave estrangeira. Em ambos os casos o banco cuida sozinho, por qualquer
caminho de escrita.

O grupo problemático é outro: RN02, RN03, RN09 e RN10 dependem de informação
que não está na linha gravada --- as demais assinaturas da conta, o plano a
duas tabelas de distância, um vínculo criado só depois do título, ou a árvore
inteira de gêneros. RN05, por fim, é de outra natureza: diz respeito ao que um
perfil pode ver, não ao que pode existir no banco. A
Seção~\ref{sec:rn-constraint} retoma esse grupo e explica o que decidimos
fazer com cada caso.

% ---------------------------------------------------------------------------
% RN14, RN15 e RN16 dependem das entidades OPCIONAIS `avaliacao` e
% `item_lista`, hoje FORA do escopo (13 entidades). Se o grupo decidir
% inclui-las para chegar a 15, inserir as linhas abaixo entre RN13 e RN17:
%
% RN14 & Um perfil avalia um título no máximo uma vez, com nota de 1 a 5. \\
% RN15 & Só pode avaliar um título quem já possui visualização registrada dele. \\
% RN16 & Um título não se repete na lista de um mesmo perfil. \\
%
% Nao deixar regra orfa referenciando entidade inexistente.
% ---------------------------------------------------------------------------

\section{Diagrama Entidade-Relacionamento}\label{sec:der}

\begin{figure}[ht]
\centering
% Substituir pelo DER exportado da ferramenta de modelagem:
% \includegraphics[width=\textwidth]{der-nextmedia.png}
\fbox{\parbox[c][0.28\textheight][c]{0.9\textwidth}{\centering\itshape
Espaço reservado para o DER da NextMedia\\[2mm]
\small (exportar para \bd{images/der-nextmedia.png})}}
\caption{Diagrama Entidade-Relacionamento da NextMedia.}
\label{fig:der}
\end{figure}

% ============================================================================
\chapter{Modelo Lógico}\label{cap:logico}
% ============================================================================

\orientacao{\textbf{Parte 2} da entrega. O professor pede o \emph{diagrama
lógico já normalizado, com possíveis justificativas}. O peso está nas
justificativas da Seção~\ref{sec:justificativas}: o diagrama sozinho não
sustenta a nota.}

\section{DER lógico}\label{sec:diagrama-logico}

\orientacao{Inserir o DER lógico exportado do ERD Tool do pgAdmin, com
tabelas, chaves primárias, chaves estrangeiras e cardinalidades. Explicar no
texto o que mudou em relação ao DER da Figura~\ref{fig:der} --- tabelas de
junção, colunas de chave estrangeira, chaves compartilhadas --- e por quê,
citando as regras de mapeamento aplicadas \citep{elmasri2019}.}

\begin{figure}[ht]
\centering
% \includegraphics[width=\textwidth]{logico-nextmedia.png}
\fbox{\parbox[c][0.28\textheight][c]{0.9\textwidth}{\centering\itshape
Espaço reservado para o DER lógico da NextMedia\\[2mm]
\small (exportar para \bd{images/logico-nextmedia.png})}}
\caption{DER lógico normalizado da NextMedia.}
\label{fig:logico}
\end{figure}

\section{Normalização}\label{sec:normalizacao}

\orientacao{Demonstrar que o esquema está na Terceira Forma Normal, em prosa,
sem transcrever a teoria. 1FN: ausência de atributos multivalorados e
compostos --- mostrar como gêneros e participações foram tratados. 2FN: nas
tabelas de chave composta, verificar que não há dependência parcial; o caso a
analisar é \bd{participacao}. 3FN: ausência de dependências transitivas ---
discutir \bd{assinatura} e \bd{episodio}. Fechar declarando a forma normal
atingida e se houve alguma desnormalização consciente.}

\section{Justificativas das decisões}\label{sec:justificativas}

\orientacao{Núcleo do capítulo e onde a nota se decide. Cada subseção precisa
desenvolver o argumento em prosa: problema, alternativas, escolha, custo
assumido. Não basta listar.}

\subsection{Mapeamento da especialização \bd{titulo}}

\orientacao{Apresentar as três estratégias clássicas: (a) tabela única com
atributo discriminador; (b) uma tabela por subtipo, sem supertipo;
(c) supertipo mais subtipos compartilhando a chave primária. Justificar a
escolha de (c): preserva a natureza total e exclusiva da especialização
(RN06), evita a proliferação de colunas nulas de (a) e mantém simples as
consultas sobre ``todos os títulos'', que (b) obrigaria a resolver por
\bd{UNION}. Declarar o custo: uma junção adicional para montar a visão
completa de um filme ou série.}

\subsection{\bd{valor\_cobrado} em \bd{assinatura}}

\orientacao{Antecipar a objeção: à primeira vista \bd{valor\_cobrado} duplica
\bd{plano.preco\_mensal} e violaria a normalização. Argumentar que não --- o
preço do plano muda ao longo do tempo e as cobranças já realizadas precisam
permanecer corretas. É dado histórico, funcionalmente dependente da
assinatura e não do plano. Sem essa explicação explícita, a decisão é lida
como erro.}

\subsection{Atributos derivados omitidos}

\orientacao{Explicar a omissão de um atributo \emph{total de temporadas} em
\bd{serie}: derivável por \bd{COUNT} sobre \bd{temporada}, armazená-lo criaria
redundância sujeita a inconsistência. Se o grupo optar por incluí-lo por
desempenho, declarar como desnormalização consciente mantida por
\textit{trigger}.}

\subsection{Chave primária composta de \bd{participacao}}

\orientacao{Justificar a chave \mbox{(\bd{id\_titulo}, \bd{id\_pessoa},
\bd{papel})}, que implementa RN11 e permite que a mesma pessoa acumule papéis
no mesmo título. Mostrar que a tabela está em 3FN: \bd{personagem} e
\bd{ordem\_credito} dependem da chave inteira, não de parte dela.}

\subsection{Relacionamento exclusivo em \bd{visualizacao}}

\orientacao{Explicar a implementação de RN12: duas chaves estrangeiras
anuláveis (\bd{id\_filme} e \bd{id\_episodio}) mais um \bd{CHECK} garantindo
que exatamente uma esteja preenchida. Apresentar a alternativa descartada ---
um supertipo ``conteúdo exibível'' acima de \bd{filme} e \bd{episodio} ---
reconhecendo que é conceitualmente mais elegante, mas mais pesada para o
escopo da disciplina.}

\subsection{Auto-relacionamentos}

\orientacao{Tratar \bd{genero.id\_genero\_pai} e
\bd{titulo.id\_titulo\_origem}: ambos aceitam NULL, pois raízes de hierarquia
e títulos sem obra de origem são casos legítimos. Discutir o risco de ciclos
(RN10), que o modelo relacional não impede declarativamente, e como seria
detectado --- consulta recursiva (\bd{WITH RECURSIVE}) ou \textit{trigger} de
validação.}

% ============================================================================
\chapter{Modelo Físico}\label{cap:fisico}
% ============================================================================

\orientacao{\textbf{Parte 3} da entrega. O professor pede código SQL para
PostgreSQL com três coisas: comandos DDL que criem o esquema e as tabelas com
\emph{todas} as restrições necessárias; comandos DML com a carga de dados
iniciais de exemplo; e três consultas \bd{SELECT} que exemplifiquem a
utilidade do minimundo. Este capítulo explica esse código; o script completo
vai no Apêndice~\ref{ap:sql}.}

\section{Convenções e criação do esquema}

\orientacao{Registrar SGBD e versão (PostgreSQL, via pgAdmin) e as convenções
de nomenclatura: identificadores em português, \bd{snake\_case}, tabelas no
singular, chaves primárias \bd{id\_<tabela>}, \textit{constraints} nomeadas
com prefixo (\bd{pk\_}, \bd{fk\_}, \bd{uq\_}, \bd{ck\_}). Justificar em uma
frase por que nomear \textit{constraints} importa: mensagens de erro legíveis
e migrações previsíveis. Comentar a Listagem~\ref{lst:schema}.}

\begin{lstlisting}[caption={Criação do esquema do projeto.},label={lst:schema}]
-- Esquema dedicado ao projeto, isolando os objetos da NextMedia
-- dos demais objetos do banco.
CREATE SCHEMA IF NOT EXISTS nextmedia;
SET search_path TO nextmedia;
\end{lstlisting}

\section{Tabelas e restrições}

\orientacao{Trazer as tabelas em listagens, agrupadas pelos mesmos quatro
núcleos usados no Capítulo~\ref{cap:conceitual} e respeitando a ordem de
dependência das chaves estrangeiras. Não despejar o arquivo inteiro aqui:
selecionar os trechos representativos --- um de cada núcleo --- e deixar o
script completo no Apêndice~\ref{ap:sql}. Cada listagem precede-se de um
parágrafo dizendo o que cria e por quê.

Os trechos que mais rendem discussão: a unicidade do \textit{e-mail} (RN01);
a chave primária compartilhada dos subtipos \bd{filme} e \bd{serie}, em que a
mesma coluna é chave primária e estrangeira (RN06); as chaves compostas das
entidades fracas (RN07, RN08); e o \bd{CHECK} do relacionamento exclusivo em
\bd{visualizacao} (RN12, RN13). A Listagem~\ref{lst:exemplo-tabela} mostra o
padrão de nomenclatura a seguir.}

\begin{lstlisting}[caption={Padrão de definição de tabela com \textit{constraints} nomeadas.},label={lst:exemplo-tabela}]
CREATE TABLE nextmedia.plano (
    id_plano          SERIAL,
    nome              VARCHAR(60)    NOT NULL,
    preco_mensal      NUMERIC(10,2)  NOT NULL,
    max_telas         SMALLINT       NOT NULL,
    max_perfis        SMALLINT       NOT NULL,
    qualidade_video   VARCHAR(10)    NOT NULL,

    CONSTRAINT pk_plano            PRIMARY KEY (id_plano),
    CONSTRAINT uq_plano_nome       UNIQUE (nome),
    CONSTRAINT ck_plano_preco      CHECK (preco_mensal >= 0),
    CONSTRAINT ck_plano_telas      CHECK (max_telas BETWEEN 1 AND 8),
    CONSTRAINT ck_plano_perfis     CHECK (max_perfis BETWEEN 1 AND 10)
);
\end{lstlisting}

\section{Das regras de negócio às restrições}\label{sec:rn-constraint}

\orientacao{Completar a Tabela~\ref{tab:rn-constraint} com todas as regras da
Tabela~\ref{tab:regras} e discutir no texto as que \emph{não} se resolvem
declarativamente --- RN03 e RN05 dependem do estado de outras tabelas e
exigem \textit{trigger} ou validação na aplicação. Explicar a escolha feita.

Fechar a seção justificando as ações de propagação das chaves estrangeiras:
RN17 pede \bd{ON DELETE CASCADE} de \bd{conta} para perfis e histórico;
RN18 pede que títulos não sejam apagados, o que sugere \bd{ON DELETE
RESTRICT} nas referências ao catálogo mais uma coluna de data de saída. O
contraste vale ser dito: dado de usuário é descartável, dado de catálogo é
patrimônio.}

\begin{table}[ht]
\centering
\caption{Mapeamento das regras de negócio para os mecanismos do modelo físico.}
\label{tab:rn-constraint}
\begin{tabularx}{0.95\textwidth}{p{0.08\textwidth} p{0.26\textwidth} X}
\toprule
\textbf{Regra} & \textbf{Mecanismo} & \textbf{Onde é aplicada} \\
\midrule
RN01 & \bd{UNIQUE}                  & \bd{conta.email} \\
RN06 & Chave primária compartilhada & \bd{filme}, \bd{serie} \\
RN11 & Chave primária composta      & \bd{participacao} \\
RN12 & \bd{CHECK}                   & \bd{visualizacao} \\
RN03 & \textit{Trigger} ou aplicação & A definir \\
\multicolumn{3}{@{}l}{\small\itshape Completar com as demais regras da Tabela~\ref{tab:regras}.}\\
\bottomrule
\end{tabularx}
\end{table}

\section{Carga de dados de exemplo}\label{sec:carga}

\orientacao{Os comandos DML que populam o banco. Descrever a estratégia da
carga em um parágrafo e informar o volume por tabela, para que o leitor saiba
que o esquema foi efetivamente executado e não apenas escrito.

Os dados precisam ser suficientes para que as consultas da
Seção~\ref{sec:consultas} retornem resultado interessante: sem pelo menos duas
séries com temporadas e episódios, várias visualizações espalhadas entre
perfis diferentes e algumas pessoas acumulando papéis, os \bd{SELECT}
devolvem tabelas vazias ou óbvias. Respeitar a ordem de dependência das
chaves estrangeiras na inserção.

Se os dados forem fictícios mas verossímeis, vale uma frase dizendo isso ---
títulos e nomes reais usados como exemplo não são um problema, mas convém
deixar claro que a base é ilustrativa.}

\begin{lstlisting}[caption={Exemplo de carga inicial.},label={lst:dml}]
-- Planos comercializados. A carga comeca pelas tabelas sem dependencia.
INSERT INTO nextmedia.plano (nome, preco_mensal, max_telas, max_perfis, qualidade_video)
VALUES ('Basico',        19.90, 1, 2, 'SD'),
       ('Intermediario', 39.90, 2, 4, 'HD'),
       ('Premium',       59.90, 4, 6, '4K');
\end{lstlisting}

\section{Consultas de exemplo}\label{sec:consultas}

\orientacao{As três consultas pedidas. Cada uma precisa de: um parágrafo
dizendo que pergunta do negócio ela responde, a listagem do \bd{SELECT} e o
resultado obtido. Consulta que só faz \bd{SELECT * FROM tabela} não demonstra
utilidade nenhuma; escolher perguntas que exijam junção, agregação ou as duas.

Sugestões que exercitam partes diferentes do modelo, a confirmar pelo grupo:
uma que percorra a especialização e as entidades fracas (por exemplo, séries
com o total de episódios por temporada); uma sobre o histórico de consumo
(títulos mais assistidos no período, ou tempo total assistido por perfil); e
uma que use o relacionamento N:N com atributos (pessoas que acumulam mais de
um papel no mesmo título, ou elenco completo de um título em ordem de
crédito). Uma consulta recursiva sobre a hierarquia de gêneros também
funciona, e mostra domínio de \bd{WITH RECURSIVE}.}

\subsection{Consulta 1}

\orientacao{Pergunta de negócio, listagem e resultado.}

\subsection{Consulta 2}

\orientacao{Pergunta de negócio, listagem e resultado.}

\subsection{Consulta 3}

\orientacao{Pergunta de negócio, listagem e resultado.}

% ============================================================================
\chapter{Considerações Finais}
% ============================================================================

\orientacao{Meia página a uma página, sem subdivisões. Retomar o que foi
construído nos três níveis, apontar as limitações conhecidas --- entidades
deixadas fora do escopo, regras delegadas à aplicação, risco de ciclo nos
auto-relacionamentos --- e indicar caminhos de evolução: índices para as
consultas mais frequentes, particionamento do histórico de visualizações,
inclusão das entidades de avaliação e lista.}

% ============================================================================
% APÊNDICE
% ============================================================================
\clearpage
\appendix

\chapter{Script SQL completo}\label{ap:sql}

\orientacao{O arquivo \bd{.sql} na íntegra, comentado, na ordem em que deve
ser executado: DDL (esquema, tabelas e restrições), DML (carga de dados
iniciais) e as três consultas de exemplo. Quando o script estiver pronto em
\bd{sql/nextmedia.sql}, descomentar a linha abaixo: ela insere o arquivo
automaticamente, evitando que documento e script entregue divirjam.}

% \lstinputlisting[caption={Script de criação do esquema NextMedia.},
%                  label={lst:script-completo}]{sql/nextmedia.sql}

% ============================================================================
% REFERÊNCIAS
% ============================================================================
\IfFileExists{refs.bib}{
    \clearpage
    \renewcommand{\bibname}{Referências bibliográficas}
    % \nocite{*} força a impressão de todas as entradas de refs.bib enquanto o
    % texto ainda não tem citações reais. REMOVER quando as seções estiverem
    % escritas e citando com \citet{} / \citep{}.
    \nocite{*}
    \bibliographystyle{plainnat}
    \bibliography{refs}
}{}

\end{document}
